\uuid{WEqb}
\titre{Base/Théorème de la base incomplète}
\auteur{hanani ; mzari}
\organisation{exo7}
\licence{CC BY-NC-SA 4.0}
\datecreate{2020-01-20}
\niveau{L1}
\module{Algèbre}
\chapitre{Espace vectoriel}
\sousChapitre{Base}
\difficulte{3}
\theme{Base et dimension, Base, Théorème de la base incomplète}
\exo7qcmid{lille-2.espaces-vectoriels.048}

\contenu{
\begin{question}
\qclassification{106.04}
\qtags{motcle=Base/Théorème de la base incomplète}
Soit $n$ et $p$ deux entiers  tels que $n >p \ge 1$, $E$ un espace vectoriel sur $\Rr$ de dimension $n$, et $v_1,v_2, \dots, v_p$ des vecteurs linéairement indépendants de $E$. Quelles sont les assertions vraies ?
\begin{answers}
\bad{$\{v_1,v_2, \dots, v_p\}$ est une base de $E$}.
\good{Il existe des vecteurs $u_1, \dots , u_k$ de $E$ tels que $\{v_1,v_2, \dots, v_p, u_1, \dots , u_k\}$ soit une base de $E$}.
\good{$\{v_1,v_2, \dots, v_{p-1}\}$ est une famille libre de $E$}.
\bad{$\{v_1,v_2, \dots, v_{p}\}$ est une famille génératrice de $E$}.
\end{answers}
\begin{explanations} Comme $\dim E=n$ et $n>p$, $\{v_1,v_2, \dots, v_p\}$ n'est pas une famille génératrice de $E$. Puisque cette famille est libre, d'après le théorème de la base incomplète, on peut la compléter pour avoir une base de $E$.
\vskip2mm
D'autre part, $\{v_1,v_2, \dots, v_{p-1}\}$ est libre, puisque toute famille extraite d'une famille libre est libre.
\end{explanations}
\end{question}
}
